\begin{abstract}
LLM agents are becoming a common interface for research, coding, and question
answering, yet their Thought--Action--Observation loop is often serial: the
model reasons, emits a tool call, then idles the GPU until the result returns.
This wait consumes 16--37\% of wall time in our workloads and 35--61\% in
prior reports.
Speculative tool execution can hide this wait, but existing systems need
auxiliary predictors, historical traces, or static workflow graphs, leaving a
gap for training-free, day-one deployment.
We observe that the model can be its own predictor: a probe forked at the
start of generation predicts Qwen3-32B's upcoming tool name with
74.6--99.6\% accuracy across five benchmarks.
We present \spork{} (\textbf{S}elf-s\textbf{P}eculative f\textbf{ORK}ing), a
training-free controller that dispatches the speculated tool call early,
overlapping its execution with the remaining chain-of-thought decode.
A cost model captures when speculation breaks even, and each component improves
one of its terms: a prefix-cache fork cuts probe cost, a confidence gate filters
mispredictions, and partial-token accept turns rejected probes into
speculative-decoding drafts.
On acceptance, the tool result is ready when reasoning ends; on rejection,
\spork{} falls back to serial execution with no correctness penalty.
On real-tool benchmarks, \spork{} cuts Qwen3-32B's GAIA $P_{95}$ by 18\%
(131.9$\to$108.1\,s); the mechanism holds across model sizes from 4B to 32B and
across dense and mixture-of-experts models, with task accuracy within 1\,pp of
baseline or better wherever measured.
\spork{} deploys as a thin controller over standard completion APIs (no
retraining, no auxiliary models, no offline traces) and is orthogonal to
token-level speculative decoding.
\spork{} is open source at \url{https://github.com/baihuajun24/spork}.
\end{abstract}
